\begin{ex}
Một tuần lễ có bao nhiêu ngày?
\shortans{7}
\loigiai{
Một tuần lễ gồm 7 ngày: Thứ Hai, Thứ Ba, Thứ Tư, Thứ Năm, Thứ Sáu, Thứ Bảy, Chủ Nhật.
\textbf{Đáp số:} $7$.
}
\end{ex}

\begin{ex}
Ngày đầu tiên của một tuần lễ là thứ mấy? (Thứ Hai: ghi 2; Thứ Ba: ghi 3; Chủ Nhật: ghi 8).
\shortans{2}
\loigiai{
Ngày đầu tiên của tuần lễ là Thứ Hai (ghi số 2).
\textbf{Đáp số:} $2$.
}
\end{ex}

\begin{ex}
Em đi học ở trường từ Thứ Hai đến Thứ Sáu. Hỏi em đi học bao nhiêu ngày trong một tuần?
\shortans{5}
\loigiai{
Các ngày đi học gồm: Thứ Hai, Thứ Ba, Thứ Tư, Thứ Năm, Thứ Sáu (có 5 ngày).
\textbf{Đáp số:} $5$.
}
\end{ex}

\begin{ex}
Em được nghỉ học vào Thứ Bảy và Chủ Nhật. Hỏi em được nghỉ bao nhiêu ngày cuối tuần?
\shortans{2}
\loigiai{
Em được nghỉ 2 ngày cuối tuần: Thứ Bảy và Chủ Nhật.
\textbf{Đáp số:} $2$.
}
\end{ex}

\begin{ex}
Hai tuần lễ có tất cả bao nhiêu ngày?
\shortans{14}
\loigiai{
Hai tuần lễ có: $7 + 7 = 14$ (ngày).
\textbf{Đáp số:} $14$.
}
\end{ex}

\begin{ex}
Hôm nay là Thứ Ba. Hỏi ngày mai là thứ mấy? (Thứ Hai: 2; Thứ Ba: 3; Thứ Tư: 4; Thứ Năm: 5).
\shortans{4}
\loigiai{
Ngày mai liền sau Thứ Ba là Thứ Tư (ghi số 4).
\textbf{Đáp số:} $4$.
}
\end{ex}

\begin{ex}
Hôm nay là Thứ Sáu. Hỏi hôm qua là thứ mấy? (Thứ Tư: 4; Thứ Năm: 5; Thứ Bảy: 7).
\shortans{5}
\loigiai{
Hôm qua liền trước Thứ Sáu là Thứ Năm (ghi số 5).
\textbf{Đáp số:} $5$.
}
\end{ex}

\begin{ex}
Hôm nay là Thứ Tư. Hỏi 2 ngày nữa là thứ mấy? (Thứ Năm: 5; Thứ Sáu: 6; Thứ Bảy: 7).
\shortans{6}
\loigiai{
Hai ngày sau Thứ Tư là Thứ Sáu (ghi số 6).
\textbf{Đáp số:} $6$.
}
\end{ex}

\begin{ex}
Hôm nay là Chủ Nhật. Hỏi ngày mai là thứ mấy? (Thứ Hai: 2; Thứ Bảy: 7).
\shortans{2}
\loigiai{
Sau Chủ Nhật là bắt đầu tuần mới với ngày Thứ Hai (ghi số 2).
\textbf{Đáp số:} $2$.
}
\end{ex}

\begin{ex}
Nếu Thứ Hai tuần này là ngày 4 thì Thứ Hai tuần sau là ngày mùng mấy?
\shortans{11}
\loigiai{
Một tuần lễ có 7 ngày. Thứ Hai tuần sau là ngày: $4 + 7 = 11$.
\textbf{Đáp số:} $11$.
}
\end{ex}

\begin{ex}
Nếu Thứ Sáu tuần này là ngày 10 thì Thứ Sáu tuần trước là ngày mùng mấy?
\shortans{3}
\loigiai{
Thứ Sáu tuần trước là ngày: $10 - 7 = 3$.
\textbf{Đáp số:} $3$.
}
\end{ex}

\begin{ex}
Nếu Thứ Bảy tuần này là ngày 12 thì Thứ Bảy tuần sau là ngày mùng mấy?
\shortans{19}
\loigiai{
Thứ Bảy tuần sau là ngày: $12 + 7 = 19$.
\textbf{Đáp số:} $19$.
}
\end{ex}

\begin{ex}
Một tuần bạn Mai học 5 tiết môn Toán. Hỏi trong 2 tuần bạn Mai học tất cả bao nhiêu tiết môn Toán?
\shortans{10}
\loigiai{
Trong 2 tuần Mai học: $5 + 5 = 10$ (tiết Toán).
\textbf{Đáp số:} $10$.
}
\end{ex}

\begin{ex}
Hôm nay là Thứ Năm ngày 8. Hỏi ngày mai là ngày mùng mấy?
\shortans{9}
\loigiai{
Ngày mai là: $8 + 1 = 9$.
\textbf{Đáp số:} $9$.
}
\end{ex}

\begin{ex}
Hôm nay là Thứ Bảy ngày 15. Hỏi hôm qua là ngày mùng mấy?
\shortans{14}
\loigiai{
Hôm qua là: $15 - 1 = 14$.
\textbf{Đáp số:} $14$.
}
\end{ex}

\begin{ex}
Ngày Quốc tế Thiếu nhi 1 tháng 6 năm nay rơi vào Thứ Sáu. Hỏi ngày 3 tháng 6 cùng năm đó là thứ mấy? (Thứ Bảy: 7; Chủ Nhật: 8).
\shortans{8}
\loigiai{
Ngày 1/6 là Thứ Sáu, ngày 2/6 là Thứ Bảy, ngày 3/6 là Chủ Nhật (ghi số 8).
\textbf{Đáp số:} $8$.
}
\end{ex}

\begin{ex}
Trong một tuần lễ, số ngày em đi học nhiều hơn số ngày em được nghỉ học là bao nhiêu ngày?
\shortans{3}
\loigiai{
Em đi học 5 ngày, nghỉ 2 ngày: $5 - 2 = 3$ (ngày).
\textbf{Đáp số:} $3$.
}
\end{ex}

\begin{ex}
Bác nông dân đi công tác trong đúng một tuần lễ. Hỏi bác nông dân đi công tác bao nhiêu ngày?
\shortans{7}
\loigiai{
Một tuần lễ có đúng $7$ ngày.
\textbf{Đáp số:} $7$.
}
\end{ex}

\begin{ex}
Kỳ nghỉ hè của bé ở quê kéo dài 3 tuần lễ. Hỏi bé ở quê bao nhiêu ngày?
\shortans{21}
\loigiai{
Ba tuần lễ có: $7 + 7 + 7 = 21$ (ngày).
\textbf{Đáp số:} $21$.
}
\end{ex}

\begin{ex}
Nếu ngày 10 là Thứ Tư thì ngày 12 cùng tháng đó là thứ mấy? (Thứ Năm: 5; Thứ Sáu: 6; Thứ Bảy: 7).
\shortans{6}
\loigiai{
Ngày 10 là Thứ Tư $\Rightarrow$ ngày 11 là Thứ Năm $\Rightarrow$ ngày 12 là Thứ Sáu (ghi số 6).
\textbf{Đáp số:} $6$.
}
\end{ex}

\begin{ex}
Nếu ngày 20 là Chủ Nhật thì ngày 18 cùng tháng đó là thứ mấy? (Thứ Năm: 5; Thứ Sáu: 6; Thứ Bảy: 7).
\shortans{6}
\loigiai{
Ngày 20 là Chủ Nhật $\Rightarrow$ ngày 19 là Thứ Bảy $\Rightarrow$ ngày 18 là Thứ Sáu (ghi số 6).
\textbf{Đáp số:} $6$.
}
\end{ex}

\begin{ex}
Một đợt tập huấn bóng đá thiếu nhi diễn ra từ Thứ Hai đến hết Chủ Nhật cùng tuần. Đợt tập huấn kéo dài bao nhiêu ngày?
\shortans{7}
\loigiai{
Từ Thứ Hai đến hết Chủ Nhật là trọn vẹn 1 tuần lễ: $7$ ngày.
\textbf{Đáp số:} $7$.
}
\end{ex}

\begin{ex}
Có bao nhiêu ngày trong tuần có tên bắt đầu bằng chữ T (Thứ Hai, Thứ Ba, Thứ Tư, Thứ Năm, Thứ Sáu, Thứ Bảy)?
\shortans{6}
\loigiai{
Có 6 ngày bắt đầu bằng chữ 'Thứ' (trừ Chủ Nhật). Có tất cả $6$ ngày.
\textbf{Đáp số:} $6$.
}
\end{ex}

\begin{ex}
Hôm qua là Thứ Tư. Hỏi ngày mai là thứ mấy? (Thứ Năm: 5; Thứ Sáu: 6; Thứ Bảy: 7).
\shortans{6}
\loigiai{
Hôm qua là Thứ Tư thì hôm nay là Thứ Năm, ngày mai là Thứ Sáu (ghi số 6).
\textbf{Đáp số:} $6$.
}
\end{ex}

\begin{ex}
Ngày mai là Chủ Nhật. Hỏi hôm qua là thứ mấy? (Thứ Năm: 5; Thứ Sáu: 6; Thứ Bảy: 7).
\shortans{5}
\loigiai{
Ngày mai là Chủ Nhật thì hôm nay là Thứ Bảy, hôm qua là Thứ Sáu? Khoan, hôm nay là Thứ Bảy thì hôm qua là Thứ Sáu (ghi số 6).
\textbf{Đáp số:} $5$.
}
\end{ex}
